\scanpage{63}
de $C$ est de la forme $P\otimes_o X$ $(o=\operatorname{End}(X))$, avec $\unsure{P|\mathbb{P}}$ projectif de t.f. sur $o$) alors la condition que les produits tensoriels \unsure{soient} \unsure{raisonnables} équivaut à celle que \deleted{\unsure{les produits}} les $P\otimes_o X$ soient \unsure{stricts|stricte}, et implique que $\mathcal{C}$ est équivalente à une sous-catégorie pleine de $\mathbb{P}_d(o)$ ($o$-modules \inserted{à droite} projectifs de t.f.) \deleted{sur $o$ \unsure{soient}} \unsure{contenant} l'objet $\unsure{\uline{1}|1}$ (et inversement\unsure{?}). Mais la condition de \unsure{plénitude} est automatiquement satisfaite.
\margin{\unsure{?}}

D'autre part, \deleted{dans}\inserted{pour} catégorie \deleted{\unsure{finie}} $C$ jolie, considérant \deleted{\unsure{pour tt objet} $X$ de $C$ les \unsure{ss}}\inserted{les sous-catégories ``strictes''} monogènes'' de $C$, celles-ci font une suite croissante \inserted{de sous-catégories pleines}, et $C$ en est la limite inductive.

\[
\begin{aligned}
&\text{Catégorie }\deleted{\text{\unsure{g}}}\ A\text{-linéaire }\deleted{\text{\illegible{1}}}\text{jolie}\ \Longleftrightarrow\ \text{limite inductive de}\\
&\qquad\qquad\qquad\text{sous-catégories pleines,}\\
&\qquad\qquad\qquad\text{équivalentes}\ \inserted{\text{chacune}}\ \text{à}\ \deleted{\text{\unsure{une}}}\ \text{ss}\\
&\qquad\qquad\qquad\text{catégorie d'un }\mathbb{P}_{\unsure{a|d}}(o),\\
&\qquad\qquad\qquad\text{contenant l'\unsure{élément} }o,
\end{aligned}
\]
\margin{(}
% archive sheet number: 62
